% © 2026 Ing. David Jaroš — CC BY-NC-ND 4.0
% UBT-AI-PROVENANCE-BEGIN
% schema: ubt-ai-provenance/v1
% tier: C_working
% ai_assistance: disclosed
% human_review: risk-based
% editorial_responsibility: Ing. David Jaroš
% policy: ../../AI_PROVENANCE.md
% notice: Working material; exhaustive human review is not claimed.
% UBT-AI-PROVENANCE-END

\documentclass[11pt]{article}
\usepackage[a4paper,margin=25mm]{geometry}
\usepackage{amsmath,amssymb,amsthm,mathtools}
\usepackage{booktabs}
\usepackage{hyperref}
\usepackage[most]{tcolorbox}

\newtheorem{theorem}{Theorem}
\newtheorem{proposition}{Proposition}
\newtheorem{corollary}{Corollary}
\newtheorem{definition}{Definition}
\newtheorem{remark}{Remark}
\newcommand{\B}{\mathbb B}
\newcommand{\avg}[1]{\overline{#1}}
\newcommand{\grade}[2]{\left\langle #1\right\rangle_{#2}}

\title{Counter-Propagating Compact Modes and Candidate\\
Electro-Geometric Response in UBT}
\author{David Jaro\v{s}}
\date{25 July 2026}

\begin{document}
\maketitle

\begin{tcolorbox}[colback=yellow!8,colframe=black!45,title=Proof status]
This is a non-canonical research track.  Compact-mode identities, the balanced
current/pressure result, the pure-Lorentz metric no-go, and the Gödel coframe
kinematics are established.  The action-level source, a non-zero balanced spin
channel, and a Gödel-type UBT solution remain open.
\end{tcolorbox}

\begin{abstract}
We separate two structures that were previously discussed as competing ideas.
Counter-propagating modes in the compact complex-time coordinate provide a
spectral source candidate; the biquaternionic/tetrad algebra determines whether
that candidate can enter the metric or connection.  A balanced pair has zero
net compact current but positive compact-gradient energy.  Hence it may source
a scalar or radion-like response, but it cannot select a rotation direction
without an additional covariant bivector or spin observable.  We prove that a
common local Lorentz rotation of the tetrad cannot change the metric, formulate
a falsifiable balanced-polarization test, and define the Gödel coframe only as a
kinematic target.  No closed-timelike-curve or enhanced-gravity claim is made.
\end{abstract}

\section{Canonical anchors and scope}
The current UBT geometric branch uses
\begin{equation}
 \Theta(q,\tau)\in\B=\mathbb C\otimes\mathbb H,
 \qquad \tau=t+i\psi,
\end{equation}
with covariant tetrad and metric readout
\begin{equation}
 E_\mu=\mathcal N_0^{-1/2}D_\mu\Theta,
 \qquad
 \frac12\left(E_\mu^\sharp E_\nu+E_\nu^\sharp E_\mu\right)
 =g_{\mu\nu}\mathbf1.
 \label{eq:metric-readout}
\end{equation}
The complementary oriented product is
\begin{equation}
 \Sigma_{\mu\nu}
 =\frac12\left(E_\mu^\sharp E_\nu-E_\nu^\sharp E_\mu\right).
 \label{eq:sigma}
\end{equation}
The track asks whether compact modes can generate a nontrivial source for
\eqref{eq:metric-readout} or for the connection.  It does not replace the
canonical action, and it does not identify \(\psi\) with an observed fifth
spacetime coordinate.

\section{Counter-propagating compact modes}
Let \(\psi\sim\psi+2\pi R_\psi\), \(n\in\mathbb Z\setminus\{0\}\), and
\(k=n/R_\psi\).  Consider
\begin{equation}
 \Theta(t,\psi)
 =A_+Q_+e^{-i\omega t+ik\psi}
 +A_-Q_-e^{-i\omega t-ik\psi},
 \label{eq:pair}
\end{equation}
where \(Q_\pm\in\B\) are fixed internal polarizations.  For the reduced
mode bookkeeping only, let \(X^\ddagger:=\overline{X^\sharp}\) denote the
Hermitian involution and define the positive weights
\begin{equation}
 N_\pm:=\grade{Q_\pm^\ddagger Q_\pm}{0}>0,
 \qquad
 W_\pm:=N_\pm|A_\pm|^2.
\end{equation}
This Hermitian pairing is used to count mode occupation and energy; it does not
replace the \(\sharp\)-anticommutator metric readout in
Eq.~\eqref{eq:metric-readout}.
For a periodic quantity define
\begin{equation}
 \avg{X}:=\frac{1}{2\pi R_\psi}
 \int_0^{2\pi R_\psi}X(\psi)\,d\psi.
\end{equation}
The nonzero winding makes all \(e^{\pm2ik\psi}\) interference terms vanish
under the circle average.

\begin{definition}[Compact current and gradient energy]
Within the reduced mode calculation define
\begin{align}
 j_\psi&:=\operatorname{Im}\,
 \grade{\avg{\Theta^\ddagger\partial_\psi\Theta}}{0},\\
 \mathcal P_\psi&:=
 \grade{\avg{(\partial_\psi\Theta)^\ddagger
                         (\partial_\psi\Theta)}}{0}.
\end{align}
The symbol \(\mathcal P_\psi\) is a positive compact-gradient quantity; its
conversion into a component of a covariant stress tensor depends on the action
and signature convention.
\end{definition}

\begin{theorem}[Exact averaged pair identities]
For the ansatz \eqref{eq:pair},
\begin{align}
 j_\psi&=k(W_+-W_-),
 \label{eq:j}\tag{K1}\\
 \mathcal P_\psi&=k^2(W_++W_-).
 \label{eq:p}\tag{K2}
\end{align}
For the usual quadratic time derivative,
\begin{equation}
 \grade{\avg{(\partial_t\Theta)^\ddagger(\partial_t\Theta)}}{0}
 =\omega^2(W_++W_-).
 \label{eq:et}
\end{equation}
\end{theorem}

\begin{proof}
Differentiate \eqref{eq:pair}.  The diagonal \(+n\) and \(-n\) terms give
respectively \(+ikW_+\) and \(-ikW_-\) in
\(\Theta^\sharp\partial_\psi\Theta\).  Cross terms contain
\(e^{\pm2ik\psi}\) and integrate to zero.  The squared derivative removes the
sign and gives \(k^2(W_++W_-)\).  The time derivative is identical for both
modes and gives \(\omega^2(W_++W_-)\).
\end{proof}

\begin{corollary}[Balanced standing pair]
If \(W_+=W_-=W>0\), then
\begin{equation}
 j_\psi=0,
 \qquad
 \mathcal P_\psi=2k^2W>0.
 \label{eq:balanced}
\end{equation}
Thus cancellation of compact momentum/current does not cancel compact-gradient
energy.
\end{corollary}

\section{Does the pair source \(\Theta\)?}
If \eqref{eq:pair} is an ansatz for the master field itself, it is not an
external source of \(\Theta\).  It becomes self-gravitating only through the
nonlinear dependence of the action, tetrad, connection, and metric on
\(\Theta\).  Schematically,
\begin{equation}
 \frac{\delta S[\Theta]}{\delta\Theta^\sharp}=0
 \quad\Longrightarrow\quad
 \text{mode energy and orientation enter only through the varied action}.
\end{equation}
If instead a distinct compact field \(C\) is introduced in an explicitly
non-canonical effective action,
\begin{equation}
 S[\Theta,C]=S_\Theta+S_C[g(\Theta),C]+S_{\rm mix}[\Theta,C],
\end{equation}
then variation gives an actual source term for \(\Theta\).  This route remains
open until \(S_C\) and \(S_{\rm mix}\) are derived or stated with controlled
effective-field assumptions.

\section{Algebra and grade audit}
The canonical carrier \(\B\cong Cl^+_{1,3}\) is closed under its associative
product.  Therefore products formed solely from \(F\in Cl^2\),
\(\Theta\in\B\), and an even spin connection remain in the even carrier.
A component derivative \(D_\mu\Theta\) has an external spacetime index but does
not by itself create an independent odd Clifford field.

An odd intermediate such as a vector plus trivector appears only after an
\emph{explicit} lift to the full Clifford module and multiplication by an odd
vector factor:
\begin{equation}
 aF=a\mathbin{\cdot}F+a\mathbin{\wedge}F
 \in Cl^1\oplus Cl^3.
\end{equation}
This identity is valid, but its use in UBT requires a stated embedding of the
Lorentz module into the full Clifford algebra.  It must not be presented as an
additional fundamental component of the eight-real-dimensional \(\Theta\).

For the present track the safer native orientation observable is the averaged
UBT bivector
\begin{equation}
 \avg{\Sigma}_{\mu\nu}
 :=\frac{1}{2\pi R_\psi}
 \int_0^{2\pi R_\psi}\Sigma_{\mu\nu}(\psi)\,d\psi.
 \label{eq:sigmabar}
\end{equation}
A balanced pair may have \(j_\psi=0\) while
\(\avg{\Sigma}_{\mu\nu}\neq0\), but this is an open polarization-existence
problem, not an established result.

\section{Metric-response no-go for a pure rotor}
Let \(e_\mu{}^a\) be a tetrad and
\(g_{\mu\nu}=e_\mu{}^ae_\nu{}^b\eta_{ab}\).  Under an infinitesimal common
Lorentz rotation
\begin{equation}
 \delta e_\mu{}^a=\lambda^a{}_b e_\mu{}^b,
 \qquad \lambda_{ab}=-\lambda_{ba},
\end{equation}
we obtain
\begin{align}
 \delta g_{\mu\nu}
 &=\lambda_{ab}e_\mu{}^be_\nu{}^a
   +\lambda_{ab}e_\nu{}^be_\mu{}^a=0.
 \label{eq:rotor-nogo}
\end{align}
Hence a bivector that merely rotates every tetrad leg by the same local
Lorentz transformation cannot tilt the physical light cone.  A successful
channel must generate symmetric tetrad strain, a nontrivial relative
left/right action, or an action-derived connection/torsion response.

\section{Pressure versus rotation}
The balanced result \eqref{eq:balanced} supplies no oriented current.  By
itself it can source only quantities that do not require a handed direction,
such as an energy density, lapse/radion contribution, or diagonal stress.
A rotating metric requires a surviving covariant orientation, for example
\(\avg{\Sigma}_{\mu\nu}\), a spin current, axial torsion, or an unbalanced
mode current.  Therefore the candidate chain is
\begin{equation}
 \begin{gathered}
 \text{balanced compact energy}
 +\text{nonzero covariant orientation}
 +\text{action-derived coupling}\\
 \longrightarrow
 \delta g_{0i}\ \text{or connection vorticity}
 \end{gathered}
 \label{eq:chain}
\end{equation}
Every arrow after the first term remains open.

\section{Gödel-type kinematic target}
A convenient Gödel coframe is
\begin{equation}
 \vartheta^0=a(dt+e^xdy),\qquad
 \vartheta^1=a\,dx,\qquad
 \vartheta^2=\frac{a}{\sqrt2}e^xdy,\qquad
 \vartheta^3=a\,dz,
 \label{eq:godel-coframe}
\end{equation}
with
\begin{equation}
 ds^2=-(\vartheta^0)^2+(\vartheta^1)^2
      +(\vartheta^2)^2+(\vartheta^3)^2.
\end{equation}
It has
\begin{equation}
 g_{ty}=-a^2e^x\neq0,
 \qquad
 d\vartheta^0=ae^x\,dx\wedge dy\neq0.
 \label{eq:godel-kin}
\end{equation}
Equations \eqref{eq:godel-coframe}--\eqref{eq:godel-kin} define a rotating
kinematic target.  They do not show that a compact pair produces this tetrad.
A UBT Gödel-type closure requires:
\begin{enumerate}
 \item a \(\Theta\) and admissible connection satisfying
       \(D_\mu\Theta=\sqrt{\mathcal N_0}E_\mu\) for the target coframe;
 \item the canonical Euler--Lagrange equation;
 \item action-derived source matching;
 \item a separate global causality analysis.
\end{enumerate}
Local tetrad representability is not dynamical selection, and rotating
kinematics is not a closed-timelike-curve theorem.

\section{Closed boundary and next theorem}
The exact result of this note is
\begin{equation}
 \boxed{
 W_+=W_-
 \Longrightarrow
 j_\psi=0,\quad \mathcal P_\psi>0,
 }
\end{equation}
together with the no-go \eqref{eq:rotor-nogo}.  The next decisive theorem is
therefore not a Gödel construction.  It is the balanced-polarization question:
\begin{equation}
 \boxed{
 j_\psi=0
 \quad\text{and}\quad
 \avg{\Sigma}_{\mu\nu}\stackrel{?}{\neq}0
 }
 \label{eq:next}
\end{equation}
for an admissible biquaternionic pair.  A negative answer reduces the track to
a compact-pressure/radion branch.  A positive answer permits, but does not yet
establish, an action-level rotation channel.

\end{document}
